
\begin{proposition}[All coprime $(3k,4l)$ multiplicities]\label{cdp:st:E11}
Fix admissible logarithmic-transformation data with multiplicities
$r=3k$, $s=4l$, where $k,l>0$ and $\gcd(r,s)=1$, in the family
constructed from the original elliptic-surface data. For every member
which is diffeomorphic to $S^6$, as in
\cite[Theorem labelled \texttt{t:main}]{Viaclovsky2026Log},
\[
 h^\bullet(X,T_X)=(1,1,1,0),
\]
its Kuranishi germ is a smooth reduced line, and varying the nonzero
linearization scale with those data fixed is locally miniversal.
For the original $(3,4)$ affine-twist family, the same statements
hold whenever $p\ne0$, including $|p|>1$.
\end{proposition}
\begin{proof}
Put $\mathfrak B_{r,s}=\mathbb P(r,s)$. The root morphism
$\sigma:\mathfrak B_{r,s}\to\mathfrak B_{3,4}$ over the identical
coarse $\mathbb P^1$ sends the two root sections to their $k$th and
$l$th powers. In line notation
$\sigma^*\mathcal O_{\mathfrak B_{3,4}}(1)
=\mathcal O_{\mathfrak B_{r,s}}(kl)$.
The vertical transition matrices and the action on fibre
$\mathcal O$-cohomology depend only on the linear monodromy.
The extra affine translations act trivially on both. The same
cusp comparison used in \ref{cdp:st:E9} therefore gives the actual bundles
\[
 E=G=\mathcal O\oplus\mathcal O(-kl),\qquad
 R^2\bar f_*\mathcal O_X=\mathcal O(-kl),
\]
and
\[
0\to\bar f^*E\to\mathcal T:=T_X(-\log W)
 \to\bar f^*L\to0,
\quad L=\mathcal O(r+s-rs),
\quad K_X=\bar f^*\mathcal O(kl-r-s).
\]
For example the normal-cover coordinate at the first finite point
is $u=v^k$, and differentiation changes the ordinary root tangent
line from the old one to $T_{\mathfrak B_{r,s}}$; the vertical
bundle and its cohomology are pulled back through $u=v^k$.
This specifies the comparison rather than identifying the two
horizontal tangent lines. At the cusp $W$ remains the same
normal-crossing divisor, so its logarithmic base correction is
$-p_0=-rs$ and its normal quotient is $\eta_*K_{\mathrm dP_6}$.

Locally $\sigma_*\mathcal O=\mathcal O$ and its positive direct
images vanish: the additional cyclic group has invariants
$\mathbb C\{v\}^{\mu_k}=\mathbb C\{u\}$, and averaging is
exact. Consequently the cohomology of the pulled-back lines
$\mathcal O(-kl)$ and $\mathcal O(-2kl)$ is the cohomology of
the original $\mathcal O(-1)$ and $\mathcal O(-2)$, hence zero.
Projection and Leray give
$h^\bullet(X,\bar f^*E)=(1,1,0,0)$.
For the horizontal line, its coarse direct image is
$\mathcal O_{\mathbb P^1}(-1)$: a vector field on a root disc
must have its coefficient divisible by the root parameter,
so $\rho_*T_{\mathfrak B_{r,s}}=T_{\mathbb P^1}(-p_1-p_2)$,
and the extra $-p_0$ gives the asserted degree. Moreover
$L$ and $L\otimes\mathcal O(-kl)$ have negative orbifold degree,
so have no global sections. Thus
$H^0(X,\bar f^*L)=H^1(X,\bar f^*L)=0$.
The two exact tangent sequences now give
$h^0(X,T_X)=h^1(X,T_X)=1$.

We compute degree three independently. On the smooth torus locus,
the boundary on a relative holomorphic one-form is the transpose
of the Kodaira--Spencer tensor. The latter is generically invertible
by \ref{cdp:st:E7}. Hence every ambient holomorphic one-form on an inverse-image
base disc is horizontal. It is the pullback of a base one-form.
At a finite root point the root map is a submersion, so that
one-form extends holomorphically in the root coordinate. At a
generic smooth point of $W$ its local equation is $q=z_1$, which
excludes a pole in its base coefficient. Therefore the canonical
pullback comparison gives
\[
 \bar f_*\Omega_X^1=\Omega_{\mathfrak B_{r,s}}^1.
\]
By Serre duality and the projection formula,
\[
 h^3(X,T_X)=h^0\bigl(\mathfrak B_{r,s},
                         \mathcal O(kl-2r-2s)\bigr)=0.
\]
For the last equality, a weighted monomial of that degree would
give nonnegative integers $a,b$ with
$3k(a+2)+4l(b+2)=kl$. Reduction modulo $k$ makes $b+2$ a positive
multiple of $k$, since $\gcd(k,4l)=1$. The second term alone
is then at least $4kl$, a contradiction. On a six-sphere
$\chi(T_X)=1$, giving $h^2=1$. For the original affine-twist
family, $r=3$, $s=4$, and the horizontal Leray terms are the
acyclic lines $\mathcal O(-5)$ and $\mathcal O(-6)$ on
$\mathbb P(3,4)$. Therefore $h^\bullet(\mathcal T)=(1,1,0,0)$,
and the quotient $\eta_*K_{\mathrm dP_6}$ gives
$h^\bullet(T_X)=(1,1,1,0)$ directly for every admissible affine
twist. This computation does not require the integral homology
to be that of a sphere.

The scale derivative changes the fourth torus period by a nonzero
constant multiple of the parameter derivative in \ref{cdp:st:E9}, so its
restricted torus class is nonzero. A first-order trivialization of
the compact member would turn the given fibration into a
first-order deformation of the fixed map $f$ with fixed source.
Such deformations are $H^0(X,f^*T_{\mathbb P^1})
=H^0(\mathbb P^1,T_{\mathbb P^1})$, by connected fibres and
the projection formula. Thus they are precisely infinitesimal
projective transformations of the base. The two multiple critical
values and the hexagonal critical value remain the same three
constant sections in this scale family; the local equations
$t=v^r$, $t-1=v^s$, $q=xyz$ show this also over the dual numbers.
The base vector field is therefore zero. The nonzero torus class
cannot be trivialized, proving the nonzero global scale derivative.
The one-dimensional tangent space and the analytic local-algebra
argument at the end of \ref{cdp:st:E9} prove smoothness, reducedness and
miniversality.
\end{proof}


